\chapter{矩阵乘法的高级优化}
\label{chap:advanced-matrix-multiplication}

\setcounter{footnote}{0}

矩阵乘法是科学、工程和图形应用中的重要操作。它的计算复杂度高，因而具有很高的数据
复用度，也蕴含大量优化机会。第 3 章以矩阵乘法说明如何把多维网格映射到数据，第 5 章
又以它说明如何用共享内存分块减少全局内存访问。本章将进一步深入矩阵乘法，介绍 GPU
架构上针对这一重要计算的高级优化。

\section{背景}
\label{sec:advanced-mm-background}

矩阵--矩阵乘法（简称矩阵乘法）广泛出现在科学与工程问题的数学表达中。直观起见，可先
看它的特例向量--矩阵乘法：第一个矩阵操作数只有一行，可视为向量 $v$。计算
$u=v\times B$ 就是把 $v$ 线性变换为 $u$；$u$ 的每个元素都是 $v$ 全部元素的线性
组合，矩阵 $B$ 的第 $j$ 列给出生成 $u_j$ 时的组合系数。例如在图形应用中，把空间点
的三维坐标向量乘以一个 $3\times3$ 矩阵，就能执行坐标旋转；若还需平移或缩放，则再
加一个向量或乘统一比例因子，便得到新观察点、方向与尺度下的坐标。

要对多个向量执行同一线性变换，可以把它们收集为矩阵 $A$，每行放一个向量，再计算
$D=A\times B$。结果 $D$ 的每一行就是一个变换后的向量。回到坐标变换示例，还可能要
缩放乘法结果并加上（同样可缩放的）矩阵 $C$，其中 $C$ 的每一行是要加到 $A$ 相应变换
结果上的向量。于是得到通式
\[
  D=\alpha A\times B+\beta C,
\]
它在线性代数和高性能计算中称为\textbf{通用矩阵乘法（GEMM）}。其中
$A\times B$ 通常是计算量最大的部分，本章专注于此。

矩阵可按行主序或列主序存储，具体格式也可能来自对原矩阵或其转置的操作（见第 6 章）。
本章假设两个输入矩阵均为行主序，不过所述优化也能用于其他存储格式组合。实际 GEMM 库
通常会支持这些组合。

优化是否适用、收益多大，还取决于矩阵维度。深度学习兴起前，应用中的矩阵乘法通常规模
适中，因为线性变换涉及的向量元素不多。第 19、20 章将看到，卷积神经网络的卷积层训练
和大语言模型的注意力层都能表达为大矩阵乘法。这些大规模计算使许多过去收益不高的优化
变得重要，本章的许多技术正是为大矩阵设计的。

\section{数据复用分析}
\label{sec:mm-data-reuse}

图 \ref{fig:15-1} 回顾第 5 章的分块矩阵乘法，并采用一般化的 tile 维度。每个线程块
负责 $C$ 中一个 $m\times n$ 输出 tile。块内线程协作遍历 $A$ 与 $B$ 的输入 tile，
每轮把 $A$ 的一个 $m\times k$ tile 和 $B$ 中对应的 $k\times n$ tile 从全局内存
装入共享内存，再计算这对 tile 对所有输出元素的部分贡献。输出 tile 分布保存在块内线程
的寄存器中；全部计算结束后，线程协作把它写回全局内存。

\begin{figure}[htbp]
  \centering\small
  \begin{tabular}{c@{\quad$\times$\quad}c@{\quad$\longrightarrow$\quad}c}
    \fbox{$A$：$m\times k$ tile} & \fbox{$B$：$k\times n$ tile}
      & \fbox{$C$：$m\times n$ tile}\\[-0.1em]
    共享内存 & 共享内存 & 线程寄存器
  \end{tabular}
  \caption{分块矩阵乘法回顾}
  \label{fig:15-1}
\end{figure}

第 5 章只分析了正方形 tile。一般情况下，每对输入 tile 为 $m\times n$ 个输出元素各
执行 $k$ 次乘法和 $k$ 次加法，共 $2mnk$ 次浮点操作；同时加载 $mk+kn=k(m+n)$ 个
浮点数，每个 4 B，共 $4k(m+n)$ B。因此算术强度为
\[
 \frac{2mnk}{4k(m+n)}=\frac{0.5mn}{m+n}\quad\text{FLOP/B}.
\]

较高算术强度更容易突破内存带宽瓶颈。增大 $m,n$，也就是增大输出 tile，便会提高这个
比值。例如 $m=n=32$ 时强度为 8 FLOP/B，在 H100 上仍受内存限制；增至
$m=n=128$ 后则为 32 FLOP/B，会转为受计算吞吐量限制。

其直觉是：大输出 tile 让相同输入被更多输出复用。若把图 \ref{fig:15-1} 的 $n$ 从
32 增至 128，输出 tile 在水平方向扩大 4 倍，一个线程块相当于计算原来的 4 个输出
tile。这 4 个输出都需要同一个 $A$ 输入 tile；原来由 4 个块分别重复加载，扩大后由一个
块只加载一次。同理，增大高度 $m$ 可减少 $B$ tile 在不同块间的重复加载。消除这些冗余
加载，就提高了数据复用与算术强度。

\section{用线程粗化处理更大的 tile}
\label{sec:mm-large-tiles}

使用更大的 tile，意味着每个线程必须计算多个输出元素、加载多个输入元素，所以它可以
视为线程粗化：细粒度并行的额外开销，是相同输入会由不同线程块重复加载。这类似第 8 章
对模板计算的粗化——同样通过扩大输出 tile 增强复用与算术强度。

图 \ref{fig:15-2} 中，每块计算 $bM\times bN$ 的\textbf{块级输出 tile}。当 tile
元素数多于块内线程数时，每线程负责多个输出。例如块级 tile 为 $128\times128$，线程块
为 $16\times16$，每线程可负责 $tM\times tN=8\times8$ 的\textbf{线程级输出 tile}。

\begin{figure}[htbp]
  \centering\small
  \begin{tabular}{c@{\quad$\supset$\quad}c@{\qquad}c}
    \fbox{块级 $C$ tile：$bM\times bN$}
    & \fbox{线程级 $C$ tile：$tM\times tN$}
    & \begin{tabular}{l}$A$ 输入：$tM\times bK$\\$B$ 输入：$bK\times tN$\end{tabular}
  \end{tabular}
  \caption{使用大 tile 的分块矩阵乘法及线程级 tile 层次}
  \label{fig:15-2}
\end{figure}

一般地，矩阵可在逻辑上形成 tile 层次，每层 tile 分配给相应线程层次。本例把块级 tile
交给线程块，把线程级 tile 交给线程；还可推理 warp 级、GPU 级乃至计算集群结点级 tile。
第 \ref{sec:mm-coalesced-store} 节会用到 warp 级 tile，本节先实现块级与线程级两层。

图 \ref{fig:15-3} 的内核先由所有线程求出块级输出 tile 在 $C$ 中的起始行、列
\code{bRow}、\code{bCol}，再按线程在线程块中的位置计算线程级 tile 起点
\code{tRow}、\code{tCol}。每线程声明 $tM\times tN$ 局部数组
\code{C_r}，用 \code{clear()} 清零。局部数组默认可能放在全局内存；只有所有
下标都是编译期常量时，编译器才能把它提升到寄存器，所以访问 \code{C_r} 的循环都
必须完全展开。

\begin{figure}[htbp]
  \centering
  \begin{minipage}{0.99\textwidth}
  \begin{lstlisting}[language=C++,numbers=left]
__global__ void mm_tiled_kernel(float *A, float *B, float *C,
                                unsigned int M, unsigned int N,
                                unsigned int K) {
    unsigned int bRow = blockIdx.y * bM;
    unsigned int bCol = blockIdx.x * bN;

    unsigned int tilesPerBlockX = bN / tN;
    unsigned int ty = threadIdx.x / tilesPerBlockX;
    unsigned int tx = threadIdx.x % tilesPerBlockX;
    unsigned int tRow = ty * tM;
    unsigned int tCol = tx * tN;

    float C_r[tM][tN];
    clear(C_r, tM, tN);

    for (unsigned int tile = 0; tile < (K + bK - 1) / bK; ++tile) {
        __shared__ float A_s[bM * bK];
        __shared__ float B_s[bK * bN];
        loadTile(&A[bRow * K + tile * bK], K, M - bRow,
                 K - tile * bK, &A_s[0], bK, bM, bK);
        loadTile(&B[tile * bK * N + bCol], N, K - tile * bK,
                 N - bCol, &B_s[0], bN, bK, bN);
        __syncthreads();

        mm(tM, tN, bK, &A_s[tRow * bK], bK,
           &B_s[tCol], bN, C_r);
        __syncthreads();
    }

    float *c = &C[(bRow + tRow) * N + bCol + tCol];
    unsigned int maxRow = (bRow + tRow < M) ? M - (bRow + tRow) : 0;
    unsigned int maxCol = (bCol + tCol < N) ? N - (bCol + tCol) : 0;
    writeTile(c, N, maxRow, maxCol, C_r, tM, tN);
}
  \end{lstlisting}
  \end{minipage}
  \caption{使用大 tile 的分块矩阵乘法代码}
  \label{fig:15-3}
\end{figure}

计算块级输出 tile 需要 $A$ 的 $bM\times K$ 横向面板和 $B$ 的 $K\times bN$ 纵向
面板，共分成 $K/bK$ 对输入 tile。线程块逐对加载 \code{A_s}、\code{B_s}，
屏障确认数据齐全后调用 \code{mm()} 累加贡献，再用第二个屏障确认所有线程已用完
本轮输入，随后才能覆盖共享内存。全部输入 tile 处理完后，\code{writeTile()} 把
线程级结果从寄存器写到全局内存。这些设备函数都会内联进主内核。

\begin{figure}[htbp]
  \centering
  \begin{minipage}{0.82\textwidth}
  \begin{lstlisting}[language=C++,numbers=left]
__device__ __forceinline__ void clear(float C_r[][tN],
                                      unsigned int m, unsigned int n) {
#pragma unroll
    for (unsigned int row = 0; row < m; ++row) {
#pragma unroll
        for (unsigned int col = 0; col < n; ++col)
            C_r[row][col] = 0.0f;
    }
}
  \end{lstlisting}
  \end{minipage}
  \caption{清零输出 tile 的代码}
  \label{fig:15-4}
\end{figure}

图 \ref{fig:15-4} 的双重循环遍历输出 tile 并清零。两个循环都标注
\code{#pragma unroll}，保证完全展开，让 \code{C_r} 以常量下标访问并进入寄存器。

图 \ref{fig:15-5} 的 \code{loadTile()} 接收原矩阵中 tile 的全局内存起始地址
\code{T}、主维度 \code{lda}、边界 \code{maxRow/maxCol}，以及共享内存地址
\code{T_s}、共享内存主维度 \code{ldas} 和 tile 的高、宽。\footnote{行主序中，
主维度（leading dimension）是同一列相邻元素在内存中相隔的元素数。本章的 tile 主维度
就是包含该 tile 的大矩阵宽度。}

\begin{figure}[htbp]
  \centering
  \begin{minipage}{0.98\textwidth}
  \begin{lstlisting}[language=C++,numbers=left]
__device__ __forceinline__ void loadTile(
    float *T, unsigned int lda, unsigned int maxRow, unsigned int maxCol,
    float *T_s, unsigned int ldas, unsigned int height, unsigned int width) {
    unsigned int rowsPerSubTile = NUM_THREADS_PER_BLOCK / width;
    unsigned int numSubtiles = height / rowsPerSubTile;
#pragma unroll
    for (unsigned int subTile = 0; subTile < numSubtiles; ++subTile) {
        unsigned int row = subTile * rowsPerSubTile + threadIdx.x / width;
        unsigned int col = threadIdx.x % width;
        if (row < maxRow && col < maxCol)
            T_s[row * ldas + col] = T[row * lda + col];
        else
            T_s[row * ldas + col] = 0.0f;
    }
}
  \end{lstlisting}
  \end{minipage}
  \caption{从全局内存向共享内存装载输入 tile 的代码}
  \label{fig:15-5}
\end{figure}

输入 tile 元素多于线程数时，每线程要加载多个元素。例如 tile 为 $128\times8=1024$
个元素、块中有 256 线程时，每线程加载 4 个。代码把输入 tile 逻辑划分为 4 个子 tile，
逐个装载；每轮线程计算所负责元素的行列下标，执行边界检查，越界则向共享内存填 0。
还可使用第 6 章的向量加载，让线程一条指令加载 4 个元素，留作练习。

图 \ref{fig:15-6} 的 \code{mm()} 计算当前输入 tile 对一个线程级输出 tile 的贡献。
一个 $tM\times tN$ 输出 tile 只需 $A$ 中 $tM\times bK$ 子 tile 和 $B$ 中
$bK\times tN$ 子 tile；称它们为对应线程的线程级输入 tile。$A$ tile 的横向尺寸与
$B$ tile 的纵向尺寸是块级乘法的内维度，本设计中也是线程级乘法的内维度 $bK$。

\begin{figure}[htbp]
  \centering
  \begin{minipage}{0.92\textwidth}
  \begin{lstlisting}[language=C++,numbers=left]
__device__ __forceinline__ void mm(unsigned int m, unsigned int n,
    unsigned int k, float *a, unsigned int lda, float *b,
    unsigned int ldb, float c[][tN]) {
#pragma unroll
    for (unsigned int row = 0; row < m; ++row) {
#pragma unroll
        for (unsigned int col = 0; col < n; ++col) {
#pragma unroll
            for (unsigned int i = 0; i < k; ++i)
                c[row][col] += a[row * lda + i] * b[i * ldb + col];
        }
    }
}
  \end{lstlisting}
  \end{minipage}
  \caption{使用输入 tile 执行矩阵乘法的代码}
  \label{fig:15-6}
\end{figure}

外两层循环遍历线程级输出元素并完全展开，使输出 tile 能存入寄存器；内层遍历输入 tile
相应的行和列，执行乘加。参数 \code{a}、\code{b} 指向块级共享内存 tile 中线程
所需子 tile 的起点，\code{lda/ldb} 用来跳到下一行。函数虽接收输出数组地址，内联
后所有下标都是常量，编译器仍可把数组放进寄存器。

\begin{figure}[htbp]
  \centering
  \begin{minipage}{0.92\textwidth}
  \begin{lstlisting}[language=C++,numbers=left]
__device__ __forceinline__ void writeTile(float *c, unsigned int ldc,
    unsigned int maxRow, unsigned int maxCol, float C_r[][tN],
    unsigned int m, unsigned int n) {
#pragma unroll
    for (unsigned int row = 0; row < m; ++row) {
#pragma unroll
        for (unsigned int col = 0; col < n; ++col)
            if (row < maxRow && col < maxCol)
                c[row * ldc + col] = C_r[row][col];
    }
}
  \end{lstlisting}
  \end{minipage}
  \caption{把输出 tile 写入全局内存的代码}
  \label{fig:15-7}
\end{figure}

\code{writeTile()} 以主内核算出的地址为起点，展开双重循环遍历寄存器输出 tile，经
边界检查后写入全局内存。它也可改成向量存储，让线程一条指令写 4 个元素，留作练习。
大块级与线程级 tile 已提高算术强度和总体性能，后续各节还会继续优化。

\section{输入 tile 的寄存器分块}
\label{sec:mm-register-tiling}

共享内存分块通过把会复用的数据只从全局内存加载一次，解决全局内存带宽瓶颈。一个线程
把输入元素从全局内存送到共享内存后，多个线程会再从共享内存把它加载到寄存器用于计算。
共享内存远快于全局内存，却仍慢于寄存器；全局内存瓶颈缓解后，共享内存访问延迟便可能
成为新瓶颈。

图 \ref{fig:15-2} 表明，一个线程会用同一输入值计算多个输出。在上一节实现中，同一
线程每计算一个受其贡献的输出元素，都会再次从块级共享内存 tile 加载这个输入值。可以
对输入 tile 应用寄存器分块：每个所需输入只从共享内存加载一次并放进寄存器，再从寄存器
复用于所有相应输出。

\begin{figure}[htbp]
  \centering
  \fbox{\begin{minipage}{0.9\textwidth}\centering
    共享内存：$A$ 与 $B$ 的块级输入 tile\\[0.35em]
    $\Downarrow$ 每线程每次加载 $A$ 的 $tM$ 元素条带与 $B$ 的 $tN$ 元素条带\\[0.35em]
    寄存器：$a_r[tM]$、$b_r[tN]$\\[0.35em]
    $\Downarrow$ 外积更新寄存器中的 $tM\times tN$ 输出 tile
  \end{minipage}}
  \caption{对输入 tile 使用寄存器分块的矩阵乘法}
  \label{fig:15-8}
\end{figure}

图 \ref{fig:15-8} 中，线程逐条遍历共享内存中属于自己的输入区段。每次把一对条带装进
寄存器，先计算它们对全部线程级输出的贡献，再加载下一对。

\begin{figure}[htbp]
  \centering
  \begin{minipage}{0.93\textwidth}
  \begin{lstlisting}[language=C++,numbers=left]
__device__ __forceinline__ void mm(unsigned int m, unsigned int n,
    unsigned int k, float *a, unsigned int lda, float *b,
    unsigned int ldb, float c[][tN]) {
#pragma unroll
    for (unsigned int i = 0; i < k; ++i) {
        float a_r[tM];
#pragma unroll
        for (unsigned int row = 0; row < m; ++row)
            a_r[row] = a[row * lda + i];

        float b_r[tN];
#pragma unroll
        for (unsigned int col = 0; col < n; ++col)
            b_r[col] = b[i * ldb + col];

#pragma unroll
        for (unsigned int row = 0; row < m; ++row) {
#pragma unroll
            for (unsigned int col = 0; col < n; ++col)
                c[row][col] += a_r[row] * b_r[col];
        }
    }
}
  \end{lstlisting}
  \end{minipage}
  \caption{对输入 tile 使用寄存器分块的代码}
  \label{fig:15-9}
\end{figure}

与图 \ref{fig:15-6} 相比，图 \ref{fig:15-9} 把内维度循环换到最外层。循环交换让线程
一次处理一对输入条带，并在装载下一对之前把当前条带用于全部输出。每轮声明局部数组
\code{a_r}、\code{b_r}；装载循环和计算循环都完全展开，使输入、输出 tile 都
进入寄存器。

寄存器分块减少了共享内存到寄存器的加载次数，缓解共享内存访问延迟与指令吞吐量瓶颈。
分块并不限于共享内存和寄存器，而是适用于内存层次中任何存在延迟、吞吐量差异的两级
存储。图 \ref{fig:15-8} 把共享内存中的块级 tile 与寄存器中的线程级 tile 正确结合；
超大矩阵的多 GPU、甚至多结点乘法还可把这一层次扩展到 GPU 与集群结点。

\section{输出 tile 的合并存储}
\label{sec:mm-coalesced-store}

大线程级输出 tile（如图 \ref{fig:15-2} 的 $8\times8$）有利于共享内存和寄存器分块
中的数据复用，却不利于写回。线程遍历自己的 $8\times8$ tile 时，同一 warp 中负责
相邻 tile 的线程会写入间隔 8 个元素的地址，访问不能合并。向量存储有所缓解：每线程
一次写 4 个连续元素，但只覆盖 tile 宽度的一半，不同线程的向量存储起点仍相隔 4 个
元素，依旧没有完全合并。

希望同一 warp 的线程负责相邻的 $4\times4$ 输出 tile，这样每线程可用一条向量指令
写出整行，各线程写出的 4 元素块彼此相邻，形成完全合并的事务。因此，把每线程原来的
$8\times8$ 输出 tile 改成四个分散的 $4\times4$ tile。

\begin{figure}[htbp]
  \centering\small
  \begin{tabular}{c@{$\Rightarrow$}c@{$\Rightarrow$}c}
    \fbox{块级 $128\times128$}
    & \fbox{8 个 warp：每个 $64\times32$}
    & \fbox{4 象限：每线程各取一个 $4\times4$}\\
    块级 tile & warp 级 tile & 四个物理子 tile 构成逻辑 $8\times8$
  \end{tabular}
  \caption{为合并存储而重新排列线程输出 tile}
  \label{fig:15-10}
\end{figure}

如图 \ref{fig:15-10}，256 线程即 8 个 warp，可按 $2\times4$ 排列，每个 warp 取得
$64\times32$ 的 warp 级输出 tile。再把它分成四个 $32\times16$ 象限，把 warp 的
32 个线程按 $8\times4$ 排列，每线程从每个象限取得一个 $4\times4$ tile。四个物理
tile 合起来仍是逻辑上的 $8\times8$ 线程级输出 tile，却能逐行执行完全合并的 4 元素
向量存储。\code{mm()} 和 \code{writeTile()} 需要相应改为遍历四个小 tile，并在
每步加载小输入条带，具体实现留作练习。

\section{消除 bank 冲突}
\label{sec:mm-bank-conflicts}

让线程负责大输出 tile 还会导致它们跨步访问共享内存，从而产生 bank 冲突。继续采用图
\ref{fig:15-10} 的 warp $8\times4$ 线程排列，前 4 个线程从共享内存中的 $A$ tile
加载相同位置。

图 \ref{fig:15-11}(a) 中，warp 0 加载第一条 $A$ 输入条带。线程 0--3 首先加载第 0
行第一个元素，线性下标 0，位于 bank 0；线程 4--7 加载第 4 行第一个元素。若共享内存
tile 宽度为 8，其下标是 $4\times8=32$，仍位于 bank $32\bmod32=0$。线程 8--11
加载第 8 行第一个元素，下标 64，也在 bank 0。八个四线程组的对应元素都落在同一 bank，
每次访问均形成 8 路冲突。

\begin{figure}[htbp]
  \centering\small
  \begin{tabular}{c|ccc}
    & 第 0 行首元素 & 第 4 行首元素 & 第 8 行首元素\\\hline
    (a) 主维度 8 & $0\to$ bank 0 & $32\to$ bank 0 & $64\to$ bank 0\\
    (b) 填充后主维度 9 & $0\to$ bank 0 & $36\to$ bank 4 & $72\to$ bank 8
  \end{tabular}
  \caption{增加填充列以消除共享内存 bank 冲突}
  \label{fig:15-11}
\end{figure}

第 6 章已介绍，可通过改变共享内存布局、为矩阵增加填充列来消除冲突。图
\ref{fig:15-11}(b) 加一列填充后，逻辑宽度仍为 8，地址计算使用的主维度却是 9。
第 0、4、8 行首元素的下标分别为 0、36、72，落入 bank 0、4、8。各组访问第二个元素
则进入 bank 1、5、9、……、29；8 路冲突由此消失。对每个四线程组访问的全部 16 个
元素作同样分析，可以验证 32 个 bank 下的冲突均被消除。

代码改动很简单，把图 \ref{fig:15-3} 的共享内存声明改为
\begin{lstlisting}[language=C++]
__shared__ float A_s[bM * (bK + 1)];
\end{lstlisting}
并把传给 \code{loadTile()} 和 \code{mm()} 的 \code{A_s} 主维度从
\code{bK} 改为 \code{bK+1}。$B_s$ 不需修改：同一 warp 加载的 $B$ 条带由
共享内存中连续的 32 个元素组成，天然分布在 32 个不同 bank 中。

\section{占用率考量}
\label{sec:mm-occupancy}

大 tile 提高算术强度，也会对 SM 资源施压，尤其是保存块级输入 tile 的共享内存，以及
保存线程级输入条带和输出 tile 的寄存器文件；二者会显著影响 SM 上的线程占用率。

共享内存压力容易缓解。第 \ref{sec:mm-data-reuse} 节推导的算术强度
$0.5mn/(m+n)$ 与 $k$ 无关，因此减小 $k$ 不会降低算术强度。例如 $m=n=128$ 时可取
$k=8$，此时 $A_s$、$B_s$ 分别为 $128\times8$ 与 $8\times128$，各占 4 KB，现代
GPU 上这点共享内存不会限制占用率。

寄存器压力则仍很高。每线程的 $8\times8$ 输出 tile 单独需要 64 个寄存器；对两个
线程级输入 tile 应用寄存器分块又各需 8 个，总计至少 80 个，尚未计入局部变量与临时值。
积极展开循环还会增加同时存活的临时值，进一步提高寄存器用量。

因此，SM 寄存器总数成为占用率的主要限制。实践中的高性能矩阵乘法内核往往会用到现代
GPU 允许的每线程最大 255 个寄存器。典型 SM 只有 64K 寄存器，于是每 SM 最多运行
256 个线程。这解释了第 \ref{sec:mm-large-tiles} 节为何选择每块 256 线程：每 SM 只能
执行一个这样的块，相对 SM 最多 2048 线程，占用率仅 12.5\%。

如此低的占用率通常并不理想，但寄存器分块对算术强度的巨大收益使它值得。低占用率有两
个主要缺点。第一，活跃线程较少，能发出的内存指令也少，可能无法充分利用内存带宽；可让
每线程使用向量加载、存储来发出更大的访问。第二，可用于容忍长延迟指令的线程变少；可
积极展开循环，消除长延迟分支并暴露更多互相独立的指令，让编译器拉开长延迟指令及其消费
者之间的距离。恰当展开和调度后，大部分机器码会成为一串独立的融合乘加，几乎无需停顿
即可占满 SM 核心。

仍须关注的长延迟操作是内存访问与屏障。算术强度高、受计算限制的内核通常能把内存延迟
与计算重叠：一个块等待输入 tile 从全局内存进入共享内存时，另一个块可执行展开后的融合
乘加。但本例每 SM 只能驻留一个块，无法这样重叠。一个块会分成两阶段：加载输入时受内存
限制，计算核心空闲；执行计算时受计算限制，内存硬件空闲。没有其他块隐藏延迟，就必须
在软件层面重叠两阶段。下一节介绍软件流水线。

\section{软件流水线}
\label{sec:mm-software-pipelining}

若单个线程块在每个 SM 上交替经历内存阶段与计算阶段，那么每一阶段都有一类资源闲置，
无法达到完整性能。软件流水线能把内存阶段与计算阶段重叠，让长延迟内存指令隐藏在计算
指令后面，更充分利用 GPU。

图 \ref{fig:15-12}(a) 展示未流水化时的指令次序：先加载输入 tile $A_0,B_0$，屏障
等待所有线程完成装载；再计算它们对输出 tile 的贡献 $C_0$；第二个屏障等所有线程用完
$A_0,B_0$，随后覆盖共享内存，加载 $A_1,B_1$。加载时浮点 ALU 空闲，计算时内存访问
硬件空闲。

\begin{figure}[htbp]
  \centering\small
  \begin{tabular}{c|l}
    (a) 原始 & 加载 $A_0B_0$；屏障；计算 $C_0$；屏障；加载 $A_1B_1$；……\\
    (b) 双缓冲 & 加载当前缓冲；计算当前缓冲与加载下一缓冲之间无假依赖屏障\\
    (c) 软件流水 & 第 $i$ 轮计算 $A_iB_i$，同时预取 $A_{i+1}B_{i+1}$
  \end{tabular}
  \caption{采用双缓冲与软件流水线前后的指令交错}
  \label{fig:15-12}
\end{figure}

理想情况下，应把全局内存加载与浮点计算指令交错，使长延迟内存指令后紧跟独立融合乘加，
让两类硬件持续忙碌。编译器之所以不能交错，是因为加载与计算之间存在屏障。加载后、使用
前的屏障维持真依赖，不能删除；使用当前 tile 后、覆盖为下一 tile 前的屏障只维持假依赖，
两项活动本可独立，只因复用同一地址才互斥。

双缓冲用不同共享内存缓冲区保存当前计算 tile 与下一轮加载 tile，消除假依赖以及相应
屏障；偶数、奇数轮交替使用两套缓冲即可。但待交错的指令仍位于不同循环迭代，编译器无法
跨轮调度。软件流水线进一步重排循环：第 $i$ 轮一边计算第 $i$ 对 tile，一边加载第
$i+1$ 对。这样编译器就能交错两类指令。

\begin{figure}[htbp]
  \centering\small
  \begin{tabular}{c|cccccc}
    时间 & 0 & 1 & 2 & 3 & 4 & 5\\\hline
    (a) 原始 & 加载 0 & 等待 & 计算 0 & 加载 1 & 等待 & 计算 1\\
    (b) 流水 & 加载 0 & \multicolumn{2}{c}{计算 0 $\parallel$ 加载 1}
      & \multicolumn{2}{c}{计算 1 $\parallel$ 加载 2} & 计算末 tile
  \end{tabular}
  \caption{软件流水线前后的时序示意}
  \label{fig:15-13}
\end{figure}

\begin{figure}[htbp]
  \centering
  \begin{minipage}{0.96\textwidth}
  \begin{lstlisting}[language=C++,numbers=left]
// Pre-fetch first iteration tiles to shared memory
loadTile(Acurr_s, A, 0, ...);
loadTile(Bcurr_s, B, 0, ...);
__syncthreads();

for (unsigned int tile = 0; tile < numTiles - 1; ++tile) {
    // Compute with current tiles
    mm(Acurr_s, Bcurr_s, C_r, ...);

    // Pre-fetch next tiles into the other buffers
    loadTile(Anext_s, A, tile + 1, ...);
    loadTile(Bnext_s, B, tile + 1, ...);
    __syncthreads();

    float (*Atmp_s)[bK] = Acurr_s;
    Acurr_s = Anext_s;
    Anext_s = Atmp_s;
    float (*Btmp_s)[bN] = Bcurr_s;
    Bcurr_s = Bnext_s;
    Bnext_s = Btmp_s;
}
// Compute with the final pair
mm(Acurr_s, Bcurr_s, C_r, ...);
  \end{lstlisting}
  \end{minipage}
  \caption{采用双缓冲与软件流水线的代码}
  \label{fig:15-14}
\end{figure}

代码先把第一对输入 tile 预取到 \code{Acurr_s/Bcurr_s} 并等待完成。每轮用当前缓冲
执行 \code{mm()}，同时把下一对预取到 \code{Anext_s/Bnext_s}；屏障确认下一对
全部装好后交换两套缓冲，旧的当前缓冲留给再下一轮复用。循环结束后单独计算最后一对。
当前计算与下一轮加载使用不同地址，二者之间不需屏障，内联后编译器可以交错
\code{mm()} 与 \code{loadTile()} 的指令，使浮点 ALU 和内存硬件同时忙碌。

软件流水线依赖编译器估计指令延迟并正确调度，但编译器对运行时条件了解有限。也可把交错
交给硬件：让部分 warp 专门访问内存、另一些 warp 专门计算，硬件多线程调度会自动交错，
称为\textbf{warp 专门化（warp specialization）}。还可使用后台在全局内存与共享内存
之间搬运数据的专用硬件，下一节继续讨论。

\section{专用软件与硬件支持}
\label{sec:mm-specialized-support}

矩阵乘法在 GPU 计算中非常普遍，硬件厂商通常提供融入本章优化的高性能库。NVIDIA 的
cuBLAS\cite{nvidia-cublas2025} 实现基础线性代数 BLAS API，包括矩阵乘法；
cuDNN\cite{chetlur2014cudnn} 提供深度神经网络例程，底层大量使用矩阵乘法；开源
CUTLASS\cite{nvidia-cutlass2025} 支持在不同层次和尺度上构造矩阵乘法，方便把它嵌入
更大的内核。较新的 cuTile 等接口允许开发者写数组式代码，由编译器利用张量性质自动应用
高级优化。实际 GPU 程序通常不从头实现矩阵乘法，但理解这些优化有助于迁移到底层相似的
其他计算。

GPU 还提供矩阵乘法专用硬件，可由编译器使用、由特殊 API 直接编程，更常见的是通过上述
库调用。自 Volta 起，部分 NVIDIA GPU 配备张量核心，每个核心用一条指令完成一次小矩阵
乘法。减少指令数能缓解译码与指令处理吞吐量瓶颈；张量核心还利用先进的 warp 级寄存器
交换硬件和较低精度浮点单元，获得极高吞吐量。

Hopper 之前，张量核心通过 warp 矩阵乘加（WMMA）指令编程，只能用于单个 warp。Hopper
引入 warp 组矩阵乘加（WGMMA），让多个 warp 协作完成更大的张量核心乘法。WGMMA 可
直接从共享内存读取操作数，缓解寄存器压力，并能与其他指令异步执行，减少对编译器调度
长延迟指令重叠的依赖。Blackwell 又引入张量内存（TMEM），这种片上存储专门保存张量
核心操作的输入、输出，进一步降低寄存器压力。

张量核心把计算吞吐量大幅提高后，性能可能转而受共享内存和全局内存的数据搬运速度限制。
除原始带宽外，数据搬运指令的处理开销以及寄存器占用对占用率的影响也很重要。第
\ref{sec:mm-coalesced-store} 节的向量存储就是有效硬件支持：合适的输出 tile 布局使线程
用相同指令数写更多数据，降低指令开销并改善合并。但普通向量加载、存储在全局内存与共享
内存之间搬运时仍要先装进寄存器、再从寄存器存出，既需两条指令又增加寄存器用量。

较新的 NVIDIA GPU 支持 LDGSTS 指令，能让向量数据不经寄存器直接从全局内存进入共享
内存。一条指令替代两条，降低处理开销；硬件还可在后台异步执行访问，更有效地与计算重叠，
不必依赖编译器交错指令。NVCC 可以识别“全局内存向量加载到寄存器、紧接着向量存入共享
内存”的代码序列并自动生成 LDGSTS。程序员也可在内核中调用 libcu++ 的
\code{cuda::memcpy_async}；实际使用还须遵守数据对齐等要求。

从 Hopper 起，NVIDIA GPU 还提供张量内存加速器（Tensor Memory Accelerator，TMA），
可在全局内存与共享内存间异步传送完整张量，不使用中间寄存器。它可以看成把 LDGSTS 从
一维向量推广到多维张量，可通过 \code{cuda::memcpy_async} 使用，也可间接调用使用
TMA 的 cuBLAS、cuDNN 或 CUTLASS 例程。

\section{小结}
\label{sec:advanced-mm-summary}

本章把第 6 章的多项优化用于矩阵乘法，逐步得到高级高性能内核。大输出 tile 提高算术
强度，使内核从内存受限转为计算受限；线程粗化使大 tile 成为可能。解决全局内存带宽
瓶颈后，又对输入 tile 使用寄存器分块，缓解共享内存访问延迟；重新排列线程到输出数据
的映射，保证寄存器结果写回全局内存时能够合并；为一个共享内存输入 tile 增加填充列，
消除 bank 冲突。

大量使用寄存器会造成很低的占用率，本章也给出相应缓解措施：向量加载、存储使少量线程
仍能发出足够大的访问以利用内存带宽；积极展开循环暴露足够多的独立计算指令，利用浮点
ALU 吞吐量；软件流水线或 warp 专门化把访问与计算指令交错，使内存延迟与计算重叠。

这些优化已经由 cuBLAS、cuDNN、CUTLASS 等库提供，库还会使用张量核心、张量内存、
异步复制与张量内存加速器等专用硬件。程序员通常应调用这些高度优化的库，而不必从头
实现矩阵乘法；不过，矩阵乘法仍是综合展示本书优化技术的极佳案例。

\phantomsection
\section*{练习}
\addcontentsline{toc}{section}{练习}

\begin{enumerate}[label=\arabic*.,leftmargin=2.7em]
  \item 修改图 \ref{fig:15-5} 的设备函数，用向量加载替代标量加载。
  \item 修改图 \ref{fig:15-7} 的设备函数，用向量存储替代标量存储。
  \item 按第 \ref{sec:mm-coalesced-store} 节所述的线程级输出 tile 重排方式，重新实现
  本章的分块矩阵乘法内核。
\end{enumerate}

\begin{chapterbibliography}{9}
  \bibitem{nvidia-cublas2025}
  NVIDIA, cuBLAS, 2025. \url{http://docs.nvidia.com/cuda/cublas}.

  \bibitem{chetlur2014cudnn}
  S. Chetlur, C. Woolley, P. Vandermersch, J. Cohen, J. Tran, B. Catanzaro,
  E. Shelhamer, cuDNN: efficient primitives for deep learning,
  \emph{arXiv preprint arXiv:1410.0759}, 2014.

  \bibitem{nvidia-cutlass2025}
  NVIDIA, CUTLASS, 2025. \url{https://docs.nvidia.com/cutlass}.
\end{chapterbibliography}
